\section{实验设计与评价指标}
\subsection{环境、数据与正确性}
实验运行于共享实验室服务器，使用NVIDIA A800-SXM4-80GB，驱动580.178.04、CUDA Toolkit 12.9.41、g++ 13.3.0和CMake 3.28.3。CPU型号字符串为AMD EPYC 9754，系统可见拓扑为双插槽、每插槽48核、每核2线程，共192个逻辑CPU；不根据型号字符串推定实际可见核心数。实验未建立GPU独占与显式NUMA控制条件，需保留共享负载带来的局限。

清洗后Amazon图有334863个顶点、925872条边、667129个三角形；roadNet-CA有1965206个顶点、2766607条边、120676个三角形；web-Google有875713个顶点、4322051条边、13391903个三角形。所有版本采用相同清洗与定向规则。

已有小图验证包括课程D1--D4、随机图及异常输入，GPU基线还检查了线程块边界、跨步循环和重复执行。预处理优化在30种输入、每种54种配置上逐项比较CSR的offsets和neighbors，而不是只比较最终计数。正式完整流程实验每次运行还使用CPU双指针计数核验GPU结果，核验在主计时之外。完整GPU测量的三角形数均通过检查。

\subsection{计时范围与主指标}
主计时从打开输入边表之前开始，到首次GPU计数、必要归约完成且最终结果回传CPU之后结束。计入读取、清洗、编号映射、排序、度数与定向、CSR和任务元数据准备、OpenMP开销、CUDA初始化、显存分配、传输、计算、同步和归约。不计入编译、数据下载解压、额外CPU正确性核验、日志输出及程序结束后的资源释放。

基础版为原串行主机预处理加GPU边并行二分计数；优化版仅将主机预处理切换为32线程组合模式。每版每图新进程完整运行20次，同轮数据集与版本顺序按固定种子打乱；主计时未开启分项插桩。优化加速比定义为
\begin{equation}
S_{\mathrm{opt}}=\frac{\operatorname{median}_{1\le r\le20}T_{\mathrm{base},r}}
{\operatorname{median}_{1\le r\le20}T_{\mathrm{opt},r}}.
\end{equation}
这是两组中位数的比值，不是逐对比值的中位数。当前程序参数为额外预热0次、内部重复1次；每个进程实际有一次首次GPU计算和一次额外内核采样，但完整时间在后者之前已经结束。不能把内部两次计算计为两次完整流程样本。

\subsection{消融、线程选择与记录来源}
消融每配置每图运行3次，测量的是CPU预处理时间：\texttt{fast}仅替换为快速单线程解析，\texttt{parse8}仅并行解析，\texttt{map1}拆分单线程映射，\texttt{map8}仅并行映射，\texttt{sort8}仅并行排序去重，\texttt{all8}组合三项。线程扫描分为早期2/4/8轮和后续8/16/32/64轮，不混合不同轮的统计。根据后续扫描选择统一32线程进行完整GPU比较；选择过程使用了同一批三个图，不是独立调参集，因此不宣称跨数据集最优。

本文现阶段表格依据终端汇总整理，按运行命令对应源码\texttt{e5c7e5a}，仍需原始运行元数据复核。120份逐次完整流程JSON尚未归档进仓库，不能由min/median/max推算四分位数、方差或置信区间。最终稿应补齐原始记录、输入哈希和机器状态。未进行冷磁盘/缓存清空实验，也没有提交主存峰值测量。
